% LNCS: abstract 150--250 words. No p-values in the abstract.
\begin{abstract}
ReAct-style LLM agents read content returned by external tools. Instructions planted in that content,
known as indirect prompt injection (IPI), can make the agent take unauthorised actions or leak data. Agent security
benchmarks are almost entirely in English. This paper proposes \ours, a layered defence that places
enforcement in the Tool Broker, outside the model. Its rules allow sends only to destinations written by the user,
let a confidential value leave the agent only when the user wrote that value verbatim, and redact confidential values
in the final answer. A reminder in the prompt and denial feedback support the rules. The paper states when the rules
guarantee safety and shows that every remaining bypass violates one of these conditions or is parameter tampering,
which the rules do not cover. \ours{} needs no fine-tuning and no extra LLM calls. The
paper also introduces VN-AgentInject, a Vietnamese benchmark with six attack mechanisms, benign counterparts and five
language variants. \ours{} was frozen before evaluation and run once on two disjoint sealed test sets. Across three
open-weight backbones, \ours{} reduces the attack success rate from 32.9\% to 1.2\% on the first set and from 32.4\% to
3.8\% on the second, at a cost of about 5 points of benign utility. On the first set, at the same attack
success rate, it retains about five times the utility of an estimated block-all-sinks policy. On AgentDojo, \ours{}
keeps attack success low in both English and Vietnamese but lowers utility when destinations come from tool data.

\keywords{Indirect prompt injection \and LLM agents \and Agent security \and Reference monitor \and Vietnamese}
\end{abstract}
